\documentclass[11pt]{article}
\usepackage[margin=1in]{geometry}
\usepackage{amsmath,amssymb,amsthm}
\usepackage{graphicx}
\usepackage[hidelinks]{hyperref}

\newtheorem{theorem}{Theorem}
\newtheorem{lemma}{Lemma}
\newtheorem{corollary}{Corollary}

\newcommand{\dens}{\operatorname{dens}}
\newcommand{\Lip}{\operatorname{Lip}}
\newcommand{\omegaomega}{\omega_\omega}

\title{A Conditional Higher-Smoothness Upgrade for the \\ $c_0(\Gamma)$ Embedding}
\author{agent\_lane\_15}
\date{June 18, 2026}

\begin{document}
\maketitle

\section*{Status}

This packet records a conditional reduction for a question in the 2024 paper
of Hajek, Johanis, and Schlumprecht on point character and embeddings into
$c_0(\Gamma)$ \cite{HJS240100831}.  The source paper proves that, for
the $L_p$ and Banach-lattice classes in its Corollary 28, density below
$\omega_\omega$ gives a bi-Lipschitz embedding
\[
  \Phi:X\longrightarrow c_0(\Gamma)
\]
whose scalar coordinates are $C^1$ smooth.  Immediately after that corollary,
the authors write that they do not know whether, for $p\geq 2$, the component
functions can have higher smoothness.

The present packet does not claim an unconditional full answer.  It isolates
the exact smooth-approximation input that would upgrade the authors' $C^1$
conclusion to $C^m$, and verifies that this input is supplied for WCG spaces
with a $C^m$ smooth norm by Fry and Keener \cite{FryKeener09070241}.  Thus the
remaining point for human review is the formal differentiability-order upgrade
of the smooth-coordinate $c_0$ embedding theorem used in the source paper.

\section*{Source Question}

The source paper's relevant passage is on page 11.

\begin{center}
\includegraphics[width=\textwidth]{figures/open_problem_crop.png}
\end{center}

In text form, the question is:

\begin{quote}
We do not know whether in the corollary above in the case $p\geq 2$ the
component functions can be of higher smoothness.
\end{quote}

\section*{Conditional Dependency}

The dependency below is the precise point where this packet remains
conditional.

\begin{lemma}[Formal $C^m$ smooth-coordinate upgrade]\label{lem:coordinate}
Let $m\in\mathbb N\cup\{\infty\}$.  Suppose that $X$ has the following
Lipschitz $C^m$ scalar approximation property: there is a constant $A$ such
that for every real-valued $L$-Lipschitz function $f$ on $X$ and every
$\varepsilon>0$ there exists a $C^m$ smooth function $g:X\to\mathbb R$ with
\[
  |f-g|<\varepsilon\quad\hbox{on }X,
  \qquad
  \Lip(g)\leq A L .
\]
If $X$ admits a bi-Lipschitz embedding into some $c_0(\Gamma)$, then $X$
admits a bi-Lipschitz embedding into some $c_0(\Gamma')$ whose coordinate
functions are $C^m$ smooth.
\end{lemma}

This is not proved here from first principles.  It is the formal $C^m$ version
of the $C^1$ smooth-coordinate theorem invoked by Hajek--Johanis--Schlumprecht
through \cite[Theorem 7.63]{HJS240100831} and Johanis' theorem in their proof
of Corollary 28.  The construction should be order-free: the usual proof
smooths scalar Lipschitz cut-off functions used in the $c_0$ embedding; if the
approximants are $C^m$ rather than merely $C^1$, the estimates and the
$c_0$-valued support/threshold bookkeeping are unchanged.  This last sentence
is the human-check item.

\section*{Idea of the Proof}

The source paper already supplies the metric part: for the stated $L_p$ and
Banach-lattice classes with $\dens X<\omega_\omega$, $X$ embeds
bi-Lipschitzly into $c_0(\Gamma)$.  Smoothness of the coordinates is then a
separate scalar-approximation issue.  Fry and Keener prove that on a WCG
Banach space admitting a $C^m$ smooth norm, Lipschitz real functions can be
uniformly approximated by Lipschitz $C^m$ smooth functions with controlled
Lipschitz constant.  Hence, once Lemma \ref{lem:coordinate} is accepted, every
such $X$ inherits a $c_0(\Gamma')$ bi-Lipschitz embedding with $C^m$ coordinate
functions.

For subspaces of $L_p(\mu)$, $1<p<\infty$, the structural hypotheses are
available: these spaces are reflexive, hence WCG, and they inherit the usual
$C^m$ smooth norm in every order allowed by the scalar function
$t\mapsto |t|^p$.

\section*{Conditional Theorem}

\begin{theorem}\label{thm:wcg}
Assume Lemma \ref{lem:coordinate}.  Let $m\in\mathbb N\cup\{\infty\}$ and let
$X$ be a WCG Banach space which admits a $C^m$ smooth norm.  If $X$ admits a
bi-Lipschitz embedding into $c_0(\Gamma)$, then $X$ admits a bi-Lipschitz
embedding into some $c_0(\Gamma')$ whose coordinate functions are $C^m$ smooth.
\end{theorem}

\begin{proof}
By Fry--Keener \cite{FryKeener09070241}, a WCG Banach space with a $C^m$
smooth norm has the Lipschitz $C^m$ scalar approximation property appearing in
Lemma \ref{lem:coordinate}: Lipschitz real-valued functions can be uniformly
approximated by $C^m$ smooth Lipschitz functions with a Lipschitz constant
controlled by a constant depending only on the space and the approximation
scheme.  Applying Lemma \ref{lem:coordinate} to the given bi-Lipschitz
$c_0(\Gamma)$ embedding yields the desired $C^m$ smooth-coordinate embedding
into another $c_0$ space.
\end{proof}

\begin{corollary}[Conditional $L_p$ subcase]\label{cor:lp}
Assume Lemma \ref{lem:coordinate}.  Let $X$ be a subspace of $L_p(\mu)$ with
$1<p<\infty$ and $\dens X<\omega_\omega$.  Let $m\in\mathbb N$.  If $m<p$,
then $X$ admits a bi-Lipschitz embedding into some $c_0(\Gamma')$ whose
component functions are $C^m$ smooth.  If $p$ is an even integer, the same
conclusion holds for every finite $m$.
\end{corollary}

\begin{proof}
The source paper proves in Corollary 26 that such an $X$ admits a
bi-Lipschitz embedding into $c_0(\Gamma)$ when $\dens X<\omega_\omega$.
Since $1<p<\infty$, the ambient $L_p(\mu)$ is reflexive, and therefore every
closed subspace $X$ is reflexive and WCG.  For $m<p$ the usual $L_p$ norm is
$C^m$ smooth away from zero in the standard Frechet sense; when $p$ is an even
integer the power function is polynomial and the usual norm gives all finite
orders away from zero.  Restricting the norm to $X$ gives a $C^m$ smooth norm
on $X$.  Theorem \ref{thm:wcg} completes the argument.
\end{proof}

\section*{Verification Notes}

The non-conditional parts of the argument are standard and local:
\begin{itemize}
\item $L_p(\mu)$ is reflexive for $1<p<\infty$, so its closed subspaces are
WCG.
\item Smoothness of the $L_p$ norm gives the asserted $C^m$ orders, with the
usual caveat at the origin handled by the definition of a smooth norm on
$X\setminus\{0\}$.
\item Fry--Keener's arXiv abstract states the needed controlled Lipschitz
$C^p$ approximation theorem for WCG spaces admitting a $C^p$ smooth norm.
\item The source paper supplies the underlying $c_0(\Gamma)$ bi-Lipschitz
embedding for subspaces of $L_p(\mu)$ of density below $\omega_\omega$.
\end{itemize}

The only nontrivial review item is Lemma \ref{lem:coordinate}.  A direct check
should revisit the proof of the $C^1$ coordinate-smoothing theorem cited in
Corollary 28 and verify that replacing $C^1$ scalar approximants by $C^m$
approximants preserves the support/threshold construction ensuring that the
map is $c_0$-valued and bi-Lipschitz.

\section*{Limitations}

This packet does not solve the full higher-smoothness question.  It is
conditional on Lemma \ref{lem:coordinate}; it does not cover arbitrary
super-reflexive spaces, arbitrary WCG finite-cotype spaces, or all orders of
smoothness for non-even $p$.  It also does not address the two broader
problems at the end of the source paper: the general super-reflexive/WCG
embedding problem and the question whether a bi-Lipschitz embedding of $B_X$
into $Y$ forces one of $X$ into $Y$.

\section*{Novelty and Literature Check}

The run indexes were searched for arXiv:2401.00831, point-character keywords,
smooth-coordinate keywords, and higher-smoothness phrases.  No existing packet
or attempt was found.  Web searches for exact phrases from the source question
and for the source title plus ``higher smoothness'' found the source paper
itself and no later explicit answer.  Fry--Keener \cite{FryKeener09070241} is
used as a supporting smooth-approximation theorem rather than as an explicit
answer to the source question.

\section*{Human Review Recommendation}

Review Lemma \ref{lem:coordinate} first.  If the $C^1$ smooth-coordinate
embedding theorem cited by the source paper is indeed formal in the smoothness
class once controlled Lipschitz scalar approximation is available, then
Corollary \ref{cor:lp} should be promotable from conditional to a partial
answer to the higher-smoothness question.

\begin{thebibliography}{9}

\bibitem{HJS240100831}
P. Hajek, M. Johanis, and T. Schlumprecht,
\emph{Remarks on the point character of Banach spaces and non-linear
embeddings into $c_0(\Gamma)$},
arXiv:2401.00831, 2024.

\bibitem{FryKeener09070241}
R. Fry and L. Keener,
\emph{Approximation of Lipschitz functions by Lipschitz, $C^p$ smooth
functions on weakly compactly generated Banach spaces},
arXiv:0907.0241, 2009.

\end{thebibliography}

\end{document}
